% 学术报告 Beamer 模板
% 使用方法：xelatex academic-presentation.tex

\documentclass[aspectratio=169]{beamer}

% ============================================================
% 主题设置
% ============================================================
\usetheme{Madrid}
\usecolortheme{whale}

% ============================================================
% 宏包
% ============================================================
\usepackage[UTF8]{ctex}          % 中文支持
\usepackage{amsmath,amssymb}     % 数学公式
\usepackage{graphicx}            % 图片
\usepackage{booktabs}            % 专业表格
\usepackage{listings}            % 代码块
\usepackage{tikz}                % 绘图
\usetikzlibrary{shapes,arrows,positioning}
\usepackage{hyperref}            % 超链接

% ============================================================
% 代码块样式
% ============================================================
\lstset{
  basicstyle=\ttfamily\scriptsize,
  keywordstyle=\color{blue}\bfseries,
  commentstyle=\color{gray},
  stringstyle=\color{red},
  breaklines=true,
  numbers=left,
  numberstyle=\tiny\color{gray},
  frame=single,
  backgroundcolor=\color{gray!10},
  showstringspaces=false,
}

% ============================================================
% 标题信息
% ============================================================
\title[短标题]{基于深度学习的文本分类研究}
\subtitle{学术报告}
\author[作者]{作者姓名}
\institute[机构]{某某大学 \\ 计算机科学与技术学院}
\date{\today}
\titlegraphic{
  % \includegraphics[height=1cm]{logo.png}
}

% ============================================================
% 文档开始
% ============================================================
\begin{document}

% ------------------------------------------------------------
% 封面
% ------------------------------------------------------------
\begin{frame}
  \titlepage
\end{frame}

% ------------------------------------------------------------
% 目录
% ------------------------------------------------------------
\begin{frame}{目录}
  \tableofcontents
\end{frame}

% ============================================================
% 第一部分：引言
% ============================================================
\section{引言}

% ------------------------------------------------------------
\begin{frame}{研究背景}
  \begin{itemize}
    \item 文本分类是自然语言处理的基础任务
    \item 传统方法依赖人工特征工程
    \item 深度学习方法展现出强大潜力
  \end{itemize}

  \vspace{1em}
  \begin{block}{研究动机}
    如何设计更有效的深度学习模型来提升文本分类性能？
  \end{block}
\end{frame}

% ------------------------------------------------------------
\begin{frame}{研究目标}
  \begin{quote}
    提出一种基于 Transformer 的文本分类方法，提升分类准确率并降低人工特征工程成本。
  \end{quote}

  \vspace{1em}
  \textbf{主要创新点：}
  \begin{enumerate}
    \item 动态注意力机制
    \item 多尺度特征融合
    \item 对比学习损失函数
  \end{enumerate}
\end{frame}

% ------------------------------------------------------------
\begin{frame}{相关工作}
  \begin{table}
    \centering
    \small
    \begin{tabular}{lcc}
      \toprule
      方法 & 优点 & 缺点 \\
      \midrule
      TextCNN & 速度快 & 捕获长距离依赖能力弱 \\
      LSTM & 序列建模强 & 并行化困难 \\
      BERT & 预训练强大 & 计算成本高 \\
      \bottomrule
    \end{tabular}
    \caption{现有方法对比}
  \end{table}
\end{frame}

% ============================================================
% 第二部分：方法
% ============================================================
\section{研究方法}

% ------------------------------------------------------------
\begin{frame}{模型架构}
  \begin{columns}[T]
    \begin{column}{0.5\textwidth}
      \textbf{模型结构：}
      \begin{enumerate}
        \item 预训练编码器 (BERT)
        \item 动态注意力池化层
        \item 多尺度特征融合
        \item 分类头 (MLP)
      \end{enumerate}

      \vspace{1em}
      \textbf{训练策略：}
      \begin{itemize}
        \item 学习率：$2 \times 10^{-5}$
        \item Batch size：32
        \item Epochs：10
      \end{itemize}
    \end{column}
    \begin{column}{0.48\textwidth}
      % 流程图
      \begin{tikzpicture}[
        node distance=1.5cm,
        block/.style={
          rectangle, draw, fill=blue!20,
          text width=4em, text centered,
          rounded corners, minimum height=2em
        },
        line/.style={draw, -latex}
      ]
        \node [block] (input) {输入};
        \node [block, below of=input] (bert) {BERT};
        \node [block, below of=bert] (attn) {注意力};
        \node [block, below of=attn] (fusion) {融合};
        \node [block, below of=fusion] (output) {输出};

        \path [line] (input) -- (bert);
        \path [line] (bert) -- (attn);
        \path [line] (attn) -- (fusion);
        \path [line] (fusion) -- (output);
      \end{tikzpicture}
    \end{column}
  \end{columns}
\end{frame}

% ------------------------------------------------------------
\begin{frame}{损失函数}
  总损失函数：
  \begin{equation}
    \mathcal{L} = \mathcal{L}_{CE} + \lambda \mathcal{L}_{CL}
  \end{equation}

  其中对比学习损失：
  \begin{equation}
    \mathcal{L}_{CL} = -\log \frac{\exp(\text{sim}(z_i, z_j) / \tau)}{\sum_{k=1}^{2N} \mathbb{1}_{k \neq i} \exp(\text{sim}(z_i, z_k) / \tau)}
  \end{equation}

  \vspace{0.5em}
  \begin{itemize}
    \item $\mathcal{L}_{CE}$：交叉熵损失
    \item $\mathcal{L}_{CL}$：对比学习损失
    \item $\lambda$：平衡系数
    \item $\tau$：温度参数
  \end{itemize}
\end{frame}

% ============================================================
% 第三部分：实验
% ============================================================
\section{实验}

% ------------------------------------------------------------
\begin{frame}{数据集}
  \begin{table}
    \centering
    \begin{tabular}{lccc}
      \toprule
      数据集 & 训练集 & 测试集 & 类别数 \\
      \midrule
      AG News & 120K & 7.6K & 4 \\
      IMDB & 25K & 25K & 2 \\
      SST-2 & 67K & 1.8K & 2 \\
      \bottomrule
    \end{tabular}
    \caption{实验数据集统计}
  \end{table}

  \vspace{1em}
  \textbf{评估指标：}
  \begin{itemize}
    \item 准确率 (Accuracy)
    \item 精确率 (Precision)
    \item 召回率 (Recall)
    \item F1 分数
  \end{itemize}
\end{frame}

% ------------------------------------------------------------
\begin{frame}{实验结果}
  \begin{table}
    \centering
    \begin{tabular}{lccc}
      \toprule
      方法 & AG News & IMDB & SST-2 \\
      \midrule
      TextCNN & 91.2\% & 89.5\% & 90.1\% \\
      LSTM & 92.1\% & 90.3\% & 91.2\% \\
      BERT-base & 94.5\% & 93.2\% & 93.5\% \\
      \textbf{Ours} & \textbf{96.1\%} & \textbf{95.3\%} & \textbf{95.8\%} \\
      \bottomrule
    \end{tabular}
    \caption{各方法在不同数据集上的准确率}
  \end{table}

  \vspace{0.5em}
  \begin{alertblock}{关键发现}
    我们的方法在所有数据集上均取得最优结果，平均提升 1.6\%。
  \end{alertblock}
\end{frame}

% ------------------------------------------------------------
\begin{frame}{消融实验}
  \begin{table}
    \centering
    \begin{tabular}{lcc}
      \toprule
      模型变体 & AG News & 平均 \\
      \midrule
      Full Model & \textbf{96.1\%} & \textbf{95.7\%} \\
      w/o Attention & 95.2\% & 94.8\% \\
      w/o Multi-scale & 95.5\% & 95.1\% \\
      w/o Contrastive & 95.8\% & 95.3\% \\
      \bottomrule
    \end{tabular}
    \caption{消融实验结果}
  \end{table}

  \vspace{0.5em}
  \begin{itemize}
    \item 注意力机制贡献最大 (+0.9\%)
    \item 多尺度融合次之 (+0.6\%)
    \item 对比学习损失 (+0.4\%)
  \end{itemize}
\end{frame}

% ============================================================
% 第四部分：结论
% ============================================================
\section{结论}

% ------------------------------------------------------------
\begin{frame}{主要贡献}
  \begin{enumerate}
    \item 提出了动态注意力池化机制
    \begin{itemize}
      \item 自适应地关注重要文本片段
      \item 提升特征表示质量
    \end{itemize}

    \item 设计了多尺度特征融合模块
    \begin{itemize}
      \item 捕获不同粒度的语义信息
      \item 增强模型表达能力
    \end{itemize}

    \item 引入对比学习损失函数
    \begin{itemize}
      \item 学习更鲁棒的特征表示
      \item 提升泛化能力
    \end{itemize}

    \item 在三个基准数据集上取得 SOTA
  \end{enumerate}
\end{frame}

% ------------------------------------------------------------
\begin{frame}{未来工作}
  \begin{itemize}
    \item \textbf{多语言扩展}
    \begin{itemize}
      \item 支持更多语言的文本分类
      \item 探索跨语言迁移学习
    \end{itemize}

    \item \textbf{效率优化}
    \begin{itemize}
      \item 模型压缩和量化
      \item 知识蒸馏
    \end{itemize}

    \item \textbf{应用拓展}
    \begin{itemize}
      \item 情感分析
      \item 命名实体识别
      \item 问答系统
    \end{itemize}
  \end{itemize}
\end{frame}

% ------------------------------------------------------------
\begin{frame}{代码示例}
  \begin{lstlisting}[language=Python]
import torch
from transformers import BertModel

class TextClassifier(torch.nn.Module):
    def __init__(self, num_classes):
        super().__init__()
        self.bert = BertModel.from_pretrained('bert-base')
        self.attention = DynamicAttention()
        self.classifier = torch.nn.Linear(768, num_classes)

    def forward(self, input_ids, attention_mask):
        outputs = self.bert(input_ids, attention_mask)
        pooled = self.attention(outputs.last_hidden_state)
        return self.classifier(pooled)
  \end{lstlisting}
\end{frame}

% ============================================================
% 致谢
% ============================================================
\begin{frame}
  \centering
  \Huge 谢谢！

  \vspace{2em}
  \Large Q\&A

  \vspace{2em}
  \normalsize
  联系方式：\href{mailto:email@example.com}{email@example.com} \\
  GitHub：\href{https://github.com/username}{github.com/username}
\end{frame}

\end{document}
